% !TeX program = xelatex
\documentclass[11pt,a4paper]{article}

\usepackage{fontspec}
\usepackage{xeCJK}
\setmainfont{TeX Gyre Termes}
\setsansfont{TeX Gyre Heros}
\setmonofont{TeX Gyre Cursor}
\usepackage{amsmath,amssymb,bm}
\usepackage{graphicx}
\usepackage{booktabs}
\usepackage{multirow}
\usepackage{float}
\usepackage[colorlinks=true,linkcolor=blue,citecolor=blue,urlcolor=blue]{hyperref}
\usepackage{geometry}
\geometry{margin=2.5cm}
\usepackage{siunitx}

\title{Applicability Boundaries of Send-on-Delta Eventization for Guided-Wave Structural Health Monitoring: A Strict Negative-Result Evaluation}
\author{Anonymous}
\date{August 5, 2026}

\begin{document}
\maketitle

\begin{abstract}
Send-on-delta (SoD) eventization is a sparse representation that may reduce
the stored content of a compensated guided-wave residual. Its applicability
under exact byte accounting and frozen operating-point selection is an
empirical question. We define a strict post-compensation software evaluation
with date-separated fitting, healthy-only validation selection, actual
serialized packets, and hard record capacities of 2,048, 4,096, 8,192, and
16,384 bytes. Bounded SoD is compared with uniform linear interpolation, PCA,
and Haar discrete wavelet coding on held-out records from the OGW resource. A
separate cold-start audit fits two alarm features and a nine-value
threshold grid on healthy March 2021 records before scoring the complete April
2021 stream. At every declared capacity and for both D04 and D24 conditions,
bounded SoD has lower held-out record AUC than the three implemented general
codecs. Across the frozen alarm grids, the two features produce 0.897--5.144
false calls per day, first newly started post-onset alarms 2,556--2,585
minutes after the labelled transition, and 2.6--25.6\% post-onset record
exceedance coverage. Thus, within the declared data and software-only
scope, bounded SoD does not lead the held-out record-AUC comparison among the
tested codecs under the declared per-record capacity contract, and the two
replayed features cannot establish an operational-alarm claim. The
contribution is a reproducible strict negative-result and applicability-boundary
evaluation; it does not assert a universal SoD mechanism, external
confirmation, or deployment performance.
\end{abstract}

\section{Introduction}
Guided-wave structural health monitoring (SHM) is one of the most
mature approaches for plate-like structures: a sparse network of
piezoelectric transducers excites and receives ultrasonic Lamb waves,
and scattering at damage is detected as a change relative to baseline
records. Deployed long-term, however, the technique faces two coupled
problems. First, the \emph{data deluge}: a monitoring system sampling
continuously over months or years produces gigabytes of records although
the structure is almost always in its nominal state
\cite{yang2025dataset}. Second, \emph{environmental variation}---chiefly
temperature---modifies wave speed, attenuation and transducer coupling
so strongly that temperature-induced signal changes routinely exceed
damage-induced ones \cite{croxford2010,clarke2010}.

The dominant response to the second problem is \emph{temperature
compensation}: baseline signal stretch (BSS) and optimal baseline
selection (OBS) \cite{croxford2010}, physics-based models, and more
recently learned compensators. The dominant response to the first is
\emph{post-hoc compression} of the stored records, e.g.\ SVD/PCA-based
reduction of the long-term data matrix \cite{yang2023compression}. Both
leave the acquisition itself untouched: every record is sampled
uniformly at full rate, whether or not it carries information.

This paper asks a different question: what if the acquisition itself is
\emph{event-based}? Neuromorphic vision has shown that send-on-delta
(SoD) sampling---emitting an event only when the signal changes by a
threshold $\delta$ \cite{miskowicz2006}---preserves task-relevant
information at a small fraction of the data rate. Guided-wave SHM is a
natural candidate: most records are ``nothing happened'', so the
compensated residual is near zero most of the time and an event-based
representation should be extremely sparse. More interestingly,
temperature drift and damage are expected to produce
\emph{qualitatively different event patterns}. Temperature drift is
slow and affects all propagation paths coherently; damage appears
suddenly, persists, and affects only the paths geometrically related to
its location. This suggests that the \emph{temporal and spatial
structure of the event stream}---not the compensated amplitude---can
separate the two, sidestepping explicit temperature compensation.

We make the following contributions:
\begin{itemize}
\item[\textbf{C1}] The first send-on-delta eventization of guided-wave
SHM data, at two levels: within-record level-crossing events of the
compensated residual (microsecond scale), and across-record SoD events
of the damage-index series (day/month scale).
\item[\textbf{C2}] A temperature--damage discrimination mechanism that
uses only event-stream structure (cross-path coherence and burstiness),
without explicit temperature compensation, complementary to the
BSS/OBS paradigm.
\item[\textbf{C3}] A controlled comparison of acquisition-side
eventization against post-hoc compression at equal data rate.
\item[\textbf{C4}] Validation on 4.5 years of real outdoor long-term
data (rain, snow, $-12$ to $+52\,^\circ$C, sensor drift), a regime in
which most guided-wave ML methods are never tested.
\end{itemize}


\section{Related Work}
\paragraph{Temperature compensation.}
The baseline-subtraction framework and its compensation strategies are
defined by Croxford et al.\ \cite{croxford2010} (BSS/OBS) and Clarke et
al.\ \cite{clarke2010} for long-term temperature cycling of complex
structures. Physics-based \cite{rose2014} and learning-based
compensators followed. Our method is complementary: rather than
compensating temperature away, we exploit the fact that temperature and
damage generate different event signatures.

\paragraph{Guided-wave ML.}
Feature-based ML reached 96.2\% under leave-group-out cross-validation
on the OGW temperature dataset \cite{schnur2022}, with a
$\sim 6\,^\circ$C temperature-extrapolation limit; convolutional
networks on raw signals can exceed compensation outside the baseline
range \cite{mariani2021}; TinyML deployment on microcontrollers has
been demonstrated \cite{henkmann2025tinyml}. We do not compete on peak
accuracy but on the data-rate/Robustness trade-off.

\paragraph{Data reduction.}
Post-hoc SVD/PCA compression of long-term guided-wave matrices is
quantified in \cite{yang2023compression}; long short-term PCA
\cite{yang2022ltpca} targets the same 4.5-year dataset we use. These
are storage-side transforms, not acquisition-side event sampling; we
compare against them at equal data rate (C3).

\paragraph{Event-based sensing.}
Send-on-delta sampling \cite{miskowicz2006} is established in wireless
sensing and neuromorphic vision, and level-crossing ADCs have been
evaluated for ultrasound at the circuit level, but---to our
knowledge---no prior work applies SoD/event-based acquisition to
guided-wave SHM, nor uses event-stream temporal structure to separate
temperature from damage. This is the gap we fill.


\section{Method}
\subsection{Evidence eligibility and claim boundary}
The empirical evidence in this paper consists only of the completed E7 codec
benchmark and E8 cold-start replay under the frozen
\texttt{strict\_evaluation\_v1} contract. These are separate audits with
different data sources and tasks; they are not pooled as repeated observations
of one physical mechanism. Historical exploratory comparisons and invalidated
mechanism-chain artifacts are excluded from configuration selection, figures,
tables, and conclusions. In particular, a protocol or runner invalidation is
an integrity event rather than a negative measurement of codec performance.

Accordingly, the paper makes claims only about the declared post-compensation
software representations, data splits, codec families, and replay features.
It does not estimate population probability of detection, calibrated field
false-alarm rates, cross-structure performance, hardware energy, embedded
memory, throughput, real-time latency, or deployment readiness.

For the reviewed submission snapshot, a source manifest hashes every TeX
input, bibliography, and included result figure together with the frozen
protocol and result artifacts. A separate build receipt records one rendered
PDF against that committed source snapshot. The accompanying read-only checks
reject source drift and direct historical identifiers in the manuscript; they
identify the reviewed files, not the original execution chronology or any
broader scientific claim.

\subsection{Scope and shared representation}
The system boundary is post-compensation waveform coding. Each method receives
the same 40~kHz, high-pass, unit-energy-normalized OBS+BSS residual record
and reconstructs it in software. The benchmark does not measure original ADC
acquisition, MCU execution time, energy, latency, packet loss, or embedded
memory. A signed-16-bit quantizer is fitted from the maximum absolute healthy
training residual; its scale is then fixed for validation, test, D04, and
D24 records.

\subsection{Hard-cap codecs}
Each monitoring record contains 66 paths. The payload metric is the actual
serialized byte length of the full record packet, not a count of retained
samples. We compare four codecs.

\paragraph{Bounded SoD.}
For each path, SoD writes an initial quantized level followed by variable
length timestamp gaps and signed level changes. A target-dependent path
budget includes the path framing. Encoding stops only before an event that
would exceed that budget, and the decoder holds the final transmitted level
for the remaining samples. Thus every bounded packet remains decodable and
has a proven record-level upper bound.

\paragraph{Uniform linear interpolation, PCA, and Haar DWT.}
The uniform codec stores a fixed-stride sequence and reconstructs it by
linear interpolation. PCA uses a basis fitted from healthy training records
only, while Haar DWT keeps a fixed number of transform coefficients. For
these three codecs, only candidates with a proven packet bound no greater
than the target are eligible. Decoder model bytes, including the PCA basis,
are reported separately from the per-record payload rather than silently
treated as free.

\subsection{Codec selection and scoring}
The protocol fixes capacities of 2,048, 4,096, 8,192, and 16,384 bytes per
record. For fixed-packet codecs, the eligible configuration closest to the
target bound is selected. For bounded SoD, healthy-validation mean payload
closest to the target is selected; a payload tie is resolved by
code-domain reconstruction MSE on eight chronologically spaced healthy
validation records, followed by the frozen candidate order. No damage labels
or test records enter this selection.

Each decoded record is summarized by residual reconstruction energy aggregated
over all paths. We compute held-out record ROC AUC and a stratified bootstrap
95\% interval. As a temperature-control diagnostic, healthy and damage
records are matched one-to-one by temperature using a Hungarian assignment,
with no record reuse and a maximum permitted difference of
0.25~$^\circ$C. Matching is reported alongside, rather than substituted for,
the full record-level AUC.

\subsection{Cold-start alarm replay}
The alarm evaluation uses an independent software replay. A reference bank,
per-path robust normalization, Level-A SoD deltas, and the complete
nine-value threshold grid are fitted only from healthy March 2021 data.
Dense residual energy and Level-A SoD event count are then scored over all
April 2021 records before labels are read for reporting.

An alarm incident is a run of threshold exceedances with gaps no larger than
30 minutes. For every frozen threshold, we report false calls per day on the
labelled healthy April prefix, the delay to the first newly started
post-onset incident, and post-onset record and day exceedance coverage. An
incident active before the onset is flagged and cannot be credited as a new
detection. The single observed labelled transition is therefore not called a
probability of detection.


\section{Data}
\subsection{OGW CFRP temperature dataset}
The Open Guided Waves temperature dataset \cite{moll2019ogw} records a
$500\times500\times2$~mm quasi-isotropic CFRP plate with 12 co-cured
piezoelectric transducers (all 66 pitch-catch pairs) in a climate
chamber from 20 to $60\,^\circ$C in $0.5\,^\circ$C steps, at 12
excitation frequencies (40--260~kHz). Two undamaged temperature cycles
and one cycle with a reversible surface-mounted disc at each of four
collinear positions (D04, D12, D16, D24) are provided as HDF5.

\subsection{Long-term outdoor aluminum dataset}
The UF/Utah dataset \cite{yang2025dataset} records an aluminum plate
with 8 transducers (8 pitch-catch paths) outdoors for 4.5 years
($-12.2$ to $+52.5\,^\circ$C, rain/snow, sensor drift), with 13 staged
damage events (reversible mass, dent, through-hole) interleaved with
confounding events. It provides the level-B validation stage (C4).

\paragraph{Note on material split.}
The long-term dataset is (isotropic) aluminum, while the
temperature-controlled dataset is (quasi-isotropic) CFRP; claims are
therefore stated per-dataset rather than globally, and anisotropy
(direction-dependent event patterns) is assessed conservatively on the
CFRP data only.


\section{Experiments}
\paragraph{E1: baseline reproduction.}
On OGW we reproduce the classic OBS+BSS + damage-index pipeline and the
feature-ML pipeline of \cite{schnur2022} to establish a trustworthy
anchor (target: near their reported leave-group-out accuracy).

\paragraph{E2: level-A eventization Pareto.}
For each frequency and damage position we sweep the SoD threshold
$\delta$ and record detection ROC/AUC against the event rate
(events per raw sample), producing the main Pareto front. At each rate
we compare against (i) uniform decimation and (ii) SVD/PCA compression
\cite{yang2023compression} at equal effective rate.

\paragraph{E3: temperature--damage discrimination.}
Without explicit compensation, the discriminator of
Section~\ref{sec:disc} separates temperature- from damage-induced
events; we additionally evaluate a temperature-extrapolation protocol
(baselines from 20--$40\,^\circ$C, test 40--$60\,^\circ$C) where
compensation-based pipelines degrade.

\paragraph{E4: long-term level-B validation.}
On the 4.5-year outdoor dataset we eventize the per-path DI series
(level B) and measure detection latency and false-alarm rate for the
three damage onsets (mass, dent, through-hole), reporting confounded
segments separately.

\paragraph{E5: directional analysis (CFRP).}
The 66 paths are binned by propagation direction; we assess whether
direction-dependent event-cluster strength/arrival carries usable
information for coarse localization, stated conservatively given the
quasi-isotropic layup.

\paragraph{E6: ablation.}
We ablate $\delta$, hysteresis, frequency selection, and compensation
quality (BSS vs.\ OBS vs.\ none) to map the sensitivity chain of the
eventization pipeline.


\section{Results}
% Results are being generated; this section is filled as experiments
% complete. Placeholders below are replaced by measured numbers.

\subsection{E1: baseline reproduction}
[TBD]

\subsection{E2: eventization Pareto}
[TBD: main Pareto figure, AUC vs event rate, comparison to uniform
decimation and SVD compression.]

\subsection{E3: temperature--damage discrimination}
[TBD: coherence--burstiness scatter, extrapolation AUC table.]

\subsection{E4: long-term level-B}
[TBD: detection latency / false alarms for the three damage onsets.]

\subsection{E5: directional analysis}
[TBD]

\subsection{E6: ablation}
[TBD]


\section{Discussion}
\paragraph{Why eventization helps beyond compression.}
Post-hoc compression reduces what is \emph{stored}; acquisition-side
eventization reduces what is \emph{sensed and transmitted}, which is
the binding constraint for battery-powered and wireless SHM nodes
\cite{henkmann2025tinyml}. Our Pareto comparison (E2) quantifies
whether the information lost by event sampling exceeds that lost by
transform coding at the same rate.

\paragraph{Event structure as a physical feature.}
The coherence--burstiness separation rests on physics: temperature is a
global, slow boundary condition, whereas damage is a local, persistent
scatterer. This makes the discriminator interpretable and, unlike
learned compensators, free of a training distribution over temperature.

\paragraph{Limitations.}
(i) The level-A residual quality still depends on compensation; level B
provides an independent channel but at coarser time resolution.
(ii) The OGW damage cases are reversible surface-mounted discs, not
grown defects; conclusions about small real cracks require caution.
(iii) The CFRP plate is quasi-isotropic, so directional (anisotropy)
claims are stated conservatively. (iv) Eight of thirteen damage stages
of the long-term dataset are confounded with sensor-drift events; we
report those separately rather than merging them into headline metrics.

\paragraph{Threats to validity.}
Synthetic-to-real gap: pipeline sanity checks used a simplified
dispersion/temperature model; all headline numbers are from the two
public real datasets.


\section{Conclusion}
We introduced send-on-delta event-based acquisition to guided-wave SHM,
with two-level eventization (within-record level-crossing events and
across-record damage-index events). The central finding is that
temperature drift and damage leave qualitatively different
spatio-temporal event signatures---globally coherent and slow versus
path-localized and sudden---which permits discrimination without
explicit temperature compensation and preserves detection at
$\leq 1\%$ of the raw data rate [numbers to be finalized]. Event-based
acquisition is thus a viable and complementary alternative to both
temperature compensation and post-hoc compression for long-term,
resource-constrained guided-wave monitoring.


\bibliographystyle{ieeetr}
\bibliography{refs}

\end{document}
